\chapter{Rekenen met decimale getallen}\label{ch:g5:decimalops}

Decimale getallen tel je op, trek je af en vermenigvuldig je bijna net als
hele getallen --- de hele kunst is weten waar het \omterm{def:g4:decimals:point}{decimaalteken} komt. Dit
hoofdstuk behandelt het cijferen en de magie van vermenigvuldigen of delen
door $10$, $100$ en $1000$. (Twee decimale getallen met elkaar
vermenigvuldigen wacht tot \cref{ch:g6:decimals}.)

\section{Optellen en aftrekken}

\begin{method}[Cijferen met een decimaalteken]\label{met:g5:decimalops:addsub}
\begin{enumerate}
  \item Schrijf de getallen met de \emph{\omterm{def:g4:decimals:point}{decimaaltekens} onder elkaar} (dan
        staan de eenheden onder de eenheden en de tienden onder de
        tienden);
  \item vul het kortste decimale deel aan met nullen;
  \item tel op of trek af zoals bij hele getallen;
  \item laat het \omterm{def:g4:decimals:point}{decimaalteken} van de uitkomst recht naar beneden vallen.
\end{enumerate}
\end{method}

\begin{example}\label{ex:g5:decimalops:addsub}
$27.5 + 3.86$ (vul aan: $27.5 = 27.50$):
\[
\begin{array}{r}
2\,7.5\,0 \\
+\ \ \ 3.8\,6 \\
\hline
3\,1.3\,6
\end{array}
\qquad\qquad
\begin{array}{r}
1\,2.4\,0 \\
-\ \ \ 5.6\,5 \\
\hline
\ \ \,6.7\,5
\end{array}
\]
Bij de aftrekking ($12.4 - 5.65$): aanvullen, lenen zoals gewoonlijk, en
controleren door terug op te tellen: $6.75 + 5.65 = 12.4$.
\end{example}

\begin{example}[Aanvullen uit het hoofd]\label{ex:g5:decimalops:complement}
Wat moet er bij $7.6$ om $10$ te bereiken? Klim: $7.6 \to 8$ is $0.4$, en
dan $8 \to 10$ is $2$: er moet $2.4$ bij. Handig bij wisselgeld: je betaalt
$10$ voor iets van $7.60$ en krijgt $2.40$ terug.
\end{example}

\section{Tien keer groter, tien keer kleiner}

\begin{proposition}[Keer en gedeeld door 10, 100, 1000]\label{prop:g5:decimalops:shift}
Met $10$ vermenigvuldigen maakt elk cijfer tien keer zo veel waard: de
cijfers schuiven één plaats naar links --- wat eruitziet alsof het
\emph{\omterm{def:g4:decimals:point}{decimaalteken} één plaats naar rechts gaat}. Delen doet het
omgekeerde:
\[
3.47 \times 10 = 34.7, \qquad
3.47 \times 100 = 347, \qquad
52.1 \div 10 = 5.21, \qquad
6 \div 100 = 0.06 .
\]
\end{proposition}
\admitted

\begin{omfigure}
\begin{tikzpicture}[scale=0.9]
  \draw (0,0) grid[xstep=1.5, ystep=0.75] (6.0,2.25);
  \node[font=\small] at (0.75,2.6) {tientallen};
  \node[font=\small] at (2.25,2.6) {eenheden};
  \node[font=\small] at (3.75,2.6) {tienden};
  \node[font=\small] at (5.25,2.6) {honderdsten};
  \node[font=\normalsize, omExo] at (2.25,1.85) {$3$};
  \node[font=\normalsize, omExo] at (3.75,1.85) {$4$};
  \node[font=\normalsize, omExo] at (5.25,1.85) {$7$};
  \node[font=\normalsize, omDef] at (0.75,1.1) {$3$};
  \node[font=\normalsize, omDef] at (2.25,1.1) {$4$};
  \node[font=\normalsize, omDef] at (3.75,1.1) {$7$};
  \node[right, font=\small] at (6.2,1.85) {$3.47$};
  \node[right, font=\small, omDef] at (6.2,1.1)
    {$3.47 \times 10 = 34.7$};
  \draw[-Stealth, omDef, thick] (4.4,1.7) -- (3.1,1.25);
\end{tikzpicture}

{\small Met $10$ vermenigvuldigen: elk cijfer klimt één plaats naar links.
Het \omterm{def:g4:decimals:point}{decimaalteken} beweegt eigenlijk niet --- de cijfers doen dat.}
\end{omfigure}

\begin{example}[Omzetten, nog eens]\label{ex:g5:decimalops:convert}
Maateenheden liggen machten van tien van elkaar, dus zijn omzettingen
gewoon deze verschuivingen: $2.35$ m $= 235$ cm ($\times 100$); $470$ g
$= 0.47$ kg ($\div 1000$); $1.5$ L $= 150$ cL.
\end{example}

\section{Een decimaal getal met een heel getal vermenigvuldigen}

\begin{method}[Decimaal keer heel]\label{met:g5:decimalops:mult}
Reken uit alsof er geen \omterm{def:g4:decimals:point}{decimaalteken} stond, en geef de uitkomst daarna
even veel decimalen als de decimale factor:
\[
4.35 \times 6: \quad 435 \times 6 = 2\,610
\ \to\ 4.35 \times 6 = 26.10 = 26.1 .
\]
Schat om het \omterm{def:g4:decimals:point}{decimaalteken} met vertrouwen te zetten: $4.35$ ligt dicht bij
$4$, en $4 \times 6 = 24$ --- dus $26.1$, en niet $2.61$ of $261$.
\end{method}

\begin{example}[Herhaald optellen werkt nog steeds]\label{ex:g5:decimalops:mult}
$0.75 \times 4 = 0.75 + 0.75 + 0.75 + 0.75 = 3$: vier \omterm{def:g3:division:half}{kwarten} maken een
geheel, in decimale kleren. Vermenigvuldigen met een heel getal blijft
``zoveel keer''.
\end{example}

\begin{example}[Boodschappen]\label{ex:g5:decimalops:shopping}
Drie schriften van $2.45$ per stuk en een pen van $1.20$:
\begin{enumerate}
  \item schriften: $2.45 \times 3 = 7.35$;
  \item totaal: $7.35 + 1.20 = 8.55$;
  \item wisselgeld van $10$: $10 - 8.55 = 1.45$.
\end{enumerate}
\end{example}

\section{Oefeningen}

\begin{exercise}[$\star$]\label{exo:g5:decimalops:1}
Tel cijferend op: $45.7 + 8.93$; \quad $16.25 + 3.75$; \quad
$0.86 + 12.4$.
\end{exercise}

\begin{exercise}[$\star$]\label{exo:g5:decimalops:2}
Trek cijferend af en controleer: $9.4 - 2.72$; \quad $20 - 13.45$; \quad
$6.03 - 5.9$.
\end{exercise}

\begin{exercise}[$\star$]\label{exo:g5:decimalops:3}
Vul uit het hoofd aan tot $10$ (zoals in
\cref{ex:g5:decimalops:complement}): $6.2$; \quad $3.75$; \quad $9.99$.
\end{exercise}

\begin{exercise}[$\star$]\label{exo:g5:decimalops:4}
Reken uit zonder te cijferen:
\[
7.24 \times 10, \qquad
0.58 \times 100, \qquad
34.5 \div 10, \qquad
7 \div 100, \qquad
0.3 \times 1000 .
\]
\end{exercise}

\begin{exercise}[$\star$]\label{exo:g5:decimalops:5}
Zet om: $4.05$ m in cm; \quad $325$ cm in m; \quad $0.6$ kg in g; \quad
$85$ cL in L.
\end{exercise}

\begin{exercise}[$\star$]\label{exo:g5:decimalops:6}
Schat eerst en reken dan uit: $6.2 \times 4$; \quad $3.45 \times 8$; \quad
$12.5 \times 6$.
\end{exercise}

\begin{exercise}[$\star$]\label{exo:g5:decimalops:7}
Eén rondje op een baan is $0.4$ km. Hoe lang zijn $7$ rondjes? En $25$
rondjes?
\end{exercise}

\begin{exercise}[$\star$]\label{exo:g5:decimalops:8}
Lea koopt $2$ stokbroden van $1.15$ per stuk en een taart van $12.60$. Ze
betaalt met een biljet van $20$. Reken haar totaal en haar wisselgeld uit.
\end{exercise}

\begin{exercise}[$\star$]\label{exo:g5:decimalops:9}
In een fles gaat $1.5$ L. Hoeveel liter zit er in een pak van $6$ flessen?
Een gezin drinkt $0.75$ L per dag: hoeveel dagen doet het met één pak?
(Hoeveel keer past $0.75$ in $9$?)
\end{exercise}

\begin{exercise}[$\star\star$]\label{exo:g5:decimalops:10}
Milo rekent $5.6 \times 3 = 15.18$ uit, met de gedachte ``$5 \times 3 = 15$
en $6 \times 3 = 18$''. Laat met een schatting zien dat de uitkomst niet kan
kloppen, en reken het daarna goed uit.
\end{exercise}

\begin{exercise}[$\star\star$]\label{exo:g5:decimalops:11}
De vier rondjes van een atlete duurden $65.4$ s, $64.9$ s, $65.0$ s en
$66.1$ s.
\begin{enumerate}
  \item Reken de totale tijd van de vier rondjes uit.
  \item Het record voor vier rondjes is $260$ s. Heeft ze het verbroken? Met
        hoeveel?
\end{enumerate}
\end{exercise}
